\section{Ejecución en Tiempo Real}
\label{sec:runtime}

\subsection{El loop de eventos: \texttt{rclpy.spin()}}

Una vez que el nodo está construido y todos sus suscriptores, publicadores y timers registrados, la llamada \texttt{rclpy.spin(node)} entrega el control al loop de eventos de ROS~2:

\begin{lstlisting}[style=bash, caption={Comportamiento conceptual de \texttt{rclpy.spin()}.}]
# Pseudocodigo del loop de eventos
while rclpy.ok():
    esperar_evento()          # bloquea hasta que algo ocurre
    ejecutar_callback()       # procesa el evento
    # repetir
\end{lstlisting}

El proceso permanece bloqueado en este loop indefinidamente. Los callbacks y timers son la única manera en que el código del usuario vuelve a ejecutarse después del \texttt{\_\_init\_\_}.

\subsection{Suscripciones}

Una suscripción registra un callback que ROS~2 llama automáticamente cada vez que llega un mensaje en el tópico indicado:

\begin{lstlisting}[style=python, caption={Registro de una suscripción a odometría.}]
self.odom_sub = self.create_subscription(
    Odometry,           # tipo del mensaje
    'odom',             # nombre del topico
    self.odom_callback, # funcion a llamar
    10                  # queue size
)
\end{lstlisting}

El argumento \texttt{10} es el tamaño de la cola: si los mensajes llegan más rápido de lo que el callback los procesa, se acumulan hasta 10 mensajes antes de empezar a descartarlos. DDS gestiona la entrega del mensaje desde el publicador (el nodo de odometría de Gazebo) hasta este callback, incluso si están en procesos diferentes.

\subsection{Timers}

Un timer registra un callback que se ejecuta periódicamente:

\begin{lstlisting}[style=python, caption={Registro de un timer de control a 20~Hz.}]
self.timer = self.create_timer(
    0.05,               # periodo en segundos (20 Hz)
    self.control_loop   # funcion a llamar
)
\end{lstlisting}

\subsection{Modelo de concurrencia}

Por defecto, \texttt{rclpy.spin()} utiliza un ejecutor de un solo hilo (\texttt{SingleThreadedExecutor}), lo que significa que los callbacks y timers no se ejecutan en paralelo: el loop de eventos los pone en cola y los ejecuta uno a uno. Esto tiene dos consecuencias importantes. La primera es favorable: la lectura y escritura sobre variables compartidas (\texttt{self.x}, \texttt{self.ranges}) es segura sin necesidad de locks. La segunda es una limitación: un callback lento retrasa todos los demás, por lo que el loop de control debe ser eficiente. Si se necesita paralelismo real, existe \texttt{MultiThreadedExecutor}, pero requiere manejo explícito de concurrencia.

\subsection{El rol del DDS}

DDS (Data Distribution Service) es la capa de transporte subyacente de ROS~2. Es un estándar industrial de comunicación publicador-suscriptor. Cuando un nodo arranca, DDS anuncia su existencia en la red y otros nodos que suscriben al mismo tópico se conectan automáticamente, sin configuración manual de IPs o puertos. Para la transmisión, DDS serializa los mensajes al formato wire, los transmite (localmente vía memoria compartida, o en red vía UDP) y los deserializa en el receptor. Además, gestiona políticas de calidad de servicio (confiabilidad, durabilidad, historial) configurables por tópico.

El usuario de ROS~2 no interactúa con DDS directamente; \texttt{rclpy} lo abstrae completamente. Sin embargo, entender que DDS existe explica por qué los nodos se descubren automáticamente y por qué es posible comunicar nodos en máquinas diferentes con solo configurar \texttt{ROS\_DOMAIN\_ID}.
